% ============================================================
% Paper 1 - A Curvature Decomposition of the Explicit Formula
%           for the Riemann Zeta Function
%
% Author:   Ulrich Tehrani
% Date:     May 2026
% Version:  23
% Zenodo:   https://doi.org/10.5281/zenodo.19025598
% License:  CC BY 4.0
% ============================================================

\documentclass[11pt]{article}
\usepackage{cmap}
\usepackage[T1]{fontenc}
\usepackage[utf8]{inputenc}
\usepackage[a4paper,margin=1in]{geometry}
\usepackage{amsmath,amssymb,amsthm,mathtools}
\usepackage{xcolor}
\usepackage{hyperref}
\usepackage{enumitem}
\usepackage{setspace}
\usepackage{booktabs}
\usepackage{array}
\usepackage{graphicx}

\hypersetup{colorlinks=true,
  linkcolor=blue!50!black,
  urlcolor=blue!50!black,
  citecolor=blue!50!black}
\setlength{\parskip}{0.5em}
\setlength{\parindent}{0pt}
\onehalfspacing
\emergencystretch=1em

% --------------------------------------------------------
% Theorem environments (numbered within section)
% --------------------------------------------------------
\newtheorem{theorem}{Theorem}[section]
\newtheorem{proposition}[theorem]{Proposition}
\newtheorem{lemma}[theorem]{Lemma}
\newtheorem{corollary}[theorem]{Corollary}
\newtheorem{definition}[theorem]{Definition}
\newtheorem{assumption}[theorem]{Assumption}
\newtheorem{observation}[theorem]{Observation}
\newtheorem{conjecture}[theorem]{Conjecture}
\newtheorem*{remark}{Remark}
\newtheorem{openproblem}[theorem]{Open Problem}

% --------------------------------------------------------
% Paper-specific macros
% --------------------------------------------------------
\newcommand{\Hxi}{H_{\xi}}
\newcommand{\Hlocal}{H_{\mathrm{local}}}
\newcommand{\Hdual}{H_{\mathrm{dual}}}
\newcommand{\Sad}{S_{\mathrm{ad}}}

% --------------------------------------------------------
% Title
% --------------------------------------------------------
\title{\vspace{-1.5em}\bfseries
A Curvature Decomposition of the Explicit Formula\\[0.4em]
\large for the Riemann Zeta Function}
\author{Ulrich Tehrani}
\date{Research Note \textperiodcentered\ May~2026
      \textperiodcentered\ v23}

% ============================================================
\begin{document}
\maketitle
% ============================================================

\begin{abstract}
We study the second logarithmic derivative of the completed
Riemann zeta function in the form
\[
\Hxi(\sigma,t)\;:=\;\partial_{\sigma}^{2}\log
\lvert\xi(\sigma+it)\rvert^{2},
\]
and interpret it as a second-derivative specialisation of the
classical explicit formula. This yields a prime-cutoff
decomposition $\Hxi=\Hlocal+\Hdual$ into local curvature
contributions from the archimedean place and the finite Euler
factors, together with a bookkeeping remainder $\Hdual$
defined as the finite-cutoff difference $\Hxi-\Hlocal$.
The note also discusses the automorphic Rankin--Selberg
analogue and a structural comparison of the curvature field
with Li-type positivity criteria, noting that local positivity
does not imply global Li positivity.

\textbf{What is proved.}
Local positivity at every finite place: $V_p(\sigma)\ge 0$ for
all primes $p$ and all $\sigma>0$ (Observation~\ref{obs:Vp-pos}).
Together with the classical archimedean log-convexity
$\psi_1(\sigma)\ge 0$, this yields non-negativity of the
truncated local curvature $\Hlocal(\sigma,\kappa)\ge 0$
(Corollary~\ref{cor:Hlocal-pos}).
Critical divergence on the line $\sigma=\tfrac12$:
$\Hlocal(\tfrac12,\kappa)\asymp(\log\kappa)^{2}$ as
$\kappa\to\infty$ (Lemma~\ref{lem:critical-div}).
Subcritical convergence: $\Hlocal(\sigma,\kappa)$ tends to a
finite limit as $\kappa\to\infty$ for every $\sigma>\tfrac12$
(Lemma~\ref{lem:subcrit-conv}).

\textbf{What is numerical.}
At a representative off-critical point $(\sigma,t)=(0.3,0)$
with $\kappa=1300$, $\Hlocal\approx 915$ while
$\Hxi\approx 0.09$, so the bookkeeping remainder $\Hdual$
cancels the local curvature to better than $99.99\%$
(Observation~\ref{obs:cancellation}).
Direct evaluation at the rounded ordinate $t=14.13$, at
distance approximately $5\times 10^{-3}$ from the first
non-trivial zero $\gamma_1\approx 14.1347$, gives
$\Hxi(\tfrac12,t)\approx +89\,577.92$, illustrating the
singular kernel $(s-\rho)^{-2}$ that dominates near zeros
(Observation~\ref{obs:peak}).

\textbf{What remains open.}
Whether the singular test function implicit in $\Hxi$ ---
whose Mellin side produces the kernel $(s-\rho)^{-2}$ --- can
be embedded into, or approximated by, the class of admissible
test functions for which Weil positivity is equivalent to the
Riemann Hypothesis (Open Problem~\ref{op:weil-embedding}).

\textbf{Structure.}
The formal core in
\S\S\,\ref{sec:field}--\ref{sec:divergence} is non-circular
and proof-complete.
No use of the Riemann Hypothesis, the GUE conjecture, or the
Hilbert--P\'olya postulate is made anywhere in
\S\S\,\ref{sec:field}--\ref{sec:divergence}.
\end{abstract}

\vspace{2em}

% --------------------------------------------------------
% Notation and Conventions
% --------------------------------------------------------
\section*{Notation and Conventions}
\label{sec:notation}

\emph{Critical line.}
Throughout, $s=\sigma+it$ with $\sigma=\Re(s)$ and $t=\Im(s)$;
the critical line means $\sigma=\tfrac12$.
The completed zeta function is
$\xi(s)=\tfrac12 s(s-1)\pi^{-s/2}\Gamma(s/2)\zeta(s)$, so that
$\xi(s)=\xi(1-s)$ on the whole plane.

\medskip
\emph{Global parameters.}
The variables of this paper are $\sigma=\Re(s)>0$, $t\in\mathbb{R}$ (a real imaginary part), and
the prime cutoff $\kappa\ge 2$ (varied across the divergence
and convergence statements).
The cutoff $\kappa$ runs over integers when summing
$\sum_{p\le\kappa}$ with $p$ prime; equivalently, $\kappa$
indexes the prime $p_{\pi(\kappa)}$.
No further parameters enter the formal core.

\medskip
\emph{Main objects.}
The curvature field\footnote{We use ``curvature'' as shorthand
for $\partial_\sigma^2\log\lvert\xi\rvert^2$ without any
Riemannian, K\"ahler, or Arakelov content. This should be
distinguished from the Frenet planar curvature of
$t\mapsto\zeta(\sigma+it)$ studied by
Sourmelidis--Steuding~\cite{SourmelidisSteuding24}
and from the Quillen metric second variation of Selberg
zeta functions
(Fedosova--Rowlett--Zhang~\cite{FedosovaRowlettZhang20}).}
is
\[
\Hxi(\sigma,t)\;:=\;\partial_{\sigma}^{2}\log
\lvert\xi(\sigma+it)\rvert^{2},
\]
defined for $s=\sigma+it$ away from the zeros of $\xi$
(at a zero, the kernel $(s-\rho)^{-2}$ is singular;
cf.\ Observation~\ref{obs:peak}).
The truncated local curvature at cutoff $\kappa$ is
\[
\Hlocal(\sigma,\kappa)\;:=\;\psi_1(\sigma)
\;+\;\sum_{p\le\kappa}V_p(\sigma),
\]
with archimedean contribution
$\psi_1(\sigma)=\tfrac18\psi^{(1)}(\sigma/2)$
($\psi^{(1)}$ the trigamma function) and local Euler-factor
contribution $V_p(\sigma)=4(\log p)^2 p^{-2\sigma}/
(1-p^{-2\sigma})^2$.
Here ``archimedean contribution'' refers to the Gamma-factor
contribution in the local curvature model; the polynomial
factor $s(s-1)$ in the completed function $\xi(s)$ is not
included in $\Hlocal$.
The bookkeeping remainder is
$\Hdual(\sigma,t,\kappa):=\Hxi(\sigma,t)-\Hlocal(\sigma,\kappa)$.
The non-trivial zeros of $\zeta(s)$ are denoted $\rho$; the
first few imaginary ordinates are $\gamma_1\approx 14.135$,
$\gamma_2\approx 21.022$.

\medskip
\emph{Analytical foundations.}
The formal core uses the Hadamard product representation
of~$\xi$, log-convexity of the Gamma function
(Bohr--Mollerup; Artin~\cite{Artin}), the prime number
theorem via a Mertens-type partial-summation estimate
(Montgomery--Vaughan~\cite{MontgomeryVaughan}), and the
Tate local Mellin-norm framework~\cite{Tate1967} for the
Euler-factor contributions.

\medskip
\emph{Numerical computations.}
All numerical results were computed with \texttt{mpmath}
at 50~decimal digits of working precision;
Riemann zeros $\gamma_k$ were evaluated via
\texttt{mpmath.zetazero}.
Verification code is available at
\url{https://github.com/utehrani/analysislab-nt}.

\medskip
\emph{Not to be confused with.}
The truncated local curvature $\Hlocal(\sigma,\kappa)$ used in
this paper is the prime-side sum of squared local Euler-factor
contributions associated with the second-derivative
specialisation of the explicit formula; it is not the
classical von~Mangoldt prime sum
$\sum_{p^k\le x}\Lambda(p^k)$ and not the renormalised prime
weight that appears in admissible-test-function constructions
for the Weil functional.
The local positivity $V_p(\sigma)\ge 0$ established below is
a place-by-place positivity statement and does \emph{not}
imply Li positivity of the coefficients $\lambda_n$
(\S\,\ref{sec:li-comparison}).

\newpage

% --------------------------------------------------------
% §1 Introduction
% --------------------------------------------------------
\section{Introduction}
\label{sec:intro}

\subsection{Background and Motivation}
\label{sec:bg}

This is the first paper in a series developing an explicit,
finite-cutoff framework for studying the distribution of the
non-trivial zeros of the Riemann zeta function $\zeta(s)$.
The organising idea of the series is that prime-side
quantities at finite cutoff carry arithmetic information
about the zero geometry that is accessible by classical
analytic means.
The present paper introduces the central object of the
series, the prime-cutoff decomposition of the second
logarithmic derivative of the completed zeta function $\xi$.

The Hadamard product representation of the entire function
$\xi$, together with logarithmic differentiation, gives
\begin{equation}\label{eq:hadamard}
(\log\xi)''(s)
\;=\;-\sum_{\rho}\frac{1}{(s-\rho)^{2}},
\end{equation}
where the sum runs over the non-trivial zeros $\rho$ of
$\zeta$.
The real part of \eqref{eq:hadamard}, evaluated as a function
of $\sigma$ at fixed $t$, defines the curvature field $\Hxi$.
Stopple~\cite{Stopple16} studies the second logarithmic
derivative of $\zeta(s)$ as a meromorphic function and
analyses the location of its complex zeros. Our paper
restricts to the real variable $\sigma$ and decomposes
$\Hxi(\sigma,t)$ via the explicit formula; the two
studies are complementary.
On the local (prime and archimedean) side, the present paper
introduces a separately constructed quantity, the truncated
local curvature $\Hlocal(\sigma,\kappa)$, built in the Tate
framework from per-place log-Mellin-norm contributions: an
archimedean trigamma term $\psi_1(\sigma)$ and a finite
Euler-factor sum $\sum_{p\le\kappa}V_p(\sigma)$ with closed
forms specified in~\S\,\ref{sec:field}.
The bookkeeping remainder is defined by
$\Hdual(\sigma,t,\kappa)=\Hxi(\sigma,t)-\Hlocal(\sigma,\kappa)$;
in this paper it serves only as a bookkeeping device for the
decomposition, and no spectral interpretation is imposed on
it.

The identity
$\Hxi=\Hlocal+\Hdual$ holds by definition of $\Hdual$; its
interpretation as a curvature-form explicit-formula
reorganisation is formal at the level of singular test
functions, specialised to the kernel $2\Re(s-\rho)^{-2}$.
The classical reference for the explicit formula and its
specialisations is Edwards~\cite{Edwards}.
Burnol's invariant-analysis approach to the explicit
formula~\cite{Burnol} treats the local terms uniformly across
all places, and the Connes--Consani trace-formula work at the
archimedean place~\cite{ConnesConsani} establishes positivity
statements there for compactly supported test functions.
Neither of these isolates the prime-cutoff second-derivative
decomposition $\Hxi=\Hlocal+\Hdual$ used here, in which, at
every finite cutoff $\kappa$, the truncated local curvature
$\Hlocal(\sigma,\kappa)$ is the sum of the archimedean
trigamma term $\psi_1(\sigma)$ and the explicit closed-form
Euler-factor contributions $V_p(\sigma)$.
The principle of decomposing a global zeta-related
object as a sum over places via local Mellin/Tate
analysis is extensively developed, from Tate's
thesis~\cite{Tate1967} through Haran~\cite{Haran90}
and Burnol's conductor-operator
papers~\cite{Burnol} to Connes--Consani's
semi-local framework~\cite{ConnesConsani}. Our
paper decomposes the real-variable function
$\Hxi(\sigma,t)$ with finite prime cutoff, a
two-parameter picture not present in those works.
Equivalently, the decomposition $\Hxi=\Hlocal+\Hdual$ is the
classical explicit formula~\cite{Edwards} truncated to primes
$p\le\kappa$ and applied \emph{formally} --- at the level of
singular test functions and distributions --- to the test
function whose Mellin side produces the kernel $(s-\rho)^{-2}$.
Rigorous embedding of this singular test function into the
admissible class $\Sad$ of the Weil functional is the subject
of Open Problem~\ref{op:weil-embedding}.
At the level of the formal identity, $\Hlocal$ is the
truncated local side, while $\Hdual$ is the residual term
required to balance the zero-side curvature field after
subtracting this truncated local contribution.

The present construction, whose finite-place positivity is
the main result of this note, should be distinguished from
established operator-theoretic programmes for the Riemann
Hypothesis: it is not a Hilbert--P\'olya spectral realisation
in the sense of Connes~\cite{Connes}, nor a Berry--Keating
type quantum-Hamiltonian
construction~\cite{BerryKeating1999}, nor a de~Branges
structure-function model~\cite{deBranges,deBranges1986}.
The vocabulary of ``curvature'' used here refers solely to the
second-derivative origin of $V_p(\sigma)$ in
$\log\lvert\xi(\sigma+it)\rvert^{2}$; no further geometric
content is asserted.
To the best of our knowledge, the prime-cutoff curvature
formulation $\Hxi=\Hlocal+\Hdual$ does not appear explicitly
in the literature (cf.\ J.-F.~Burnol, personal communication).

\subsection{Main Results}
\label{sec:main}

The four formal statements of this paper concern the
prime-side object $\Hlocal(\sigma,\kappa)$ alone.
First, $V_p(\sigma)\ge 0$ at every prime $p$ and every
$\sigma>0$ (\S\,\ref{sec:positivity}, Observation~\ref{obs:Vp-pos}),
hence $\Hlocal(\sigma,\kappa)\ge 0$ throughout
(Corollary~\ref{cor:Hlocal-pos}).
Second, on the critical line, the truncated local curvature
diverges with order $(\log\kappa)^{2}$ as $\kappa\to\infty$
(\S\,\ref{sec:divergence}, Lemma~\ref{lem:critical-div}).
Third, above the critical line the truncated local curvature
tends to a finite limit
(Lemma~\ref{lem:subcrit-conv}).
The critical line is therefore the threshold separating
divergence from convergence of the local curvature sum.

The numerical illustrations of \S\,\ref{sec:numerical}
quantify the cancellation between $\Hlocal$ and $\Hdual$
off the critical line, and the singular peak of $\Hxi$ near
non-trivial zeros.

The Rankin--Selberg discussion of \S\,\ref{sec:rs} and the
Li-coefficient comparison of \S\,\ref{sec:li-comparison} are
informational only.
Their role is to place the curvature decomposition within the
classical analytic-number-theory landscape and to delimit
explicitly what local positivity does \emph{not} imply.

The single open problem
(\S\,\ref{sec:open}, Open Problem~\ref{op:weil-embedding})
asks whether the singular test function implicit in $\Hxi$
can be embedded into the class of admissible test functions
for which Weil positivity is equivalent to the Riemann
Hypothesis.

\subsection{Reader Contract}
\label{sec:contract}

This paper \emph{proves} four statements about the truncated
local curvature: place-by-place positivity at finite primes
(Observation~\ref{obs:Vp-pos}); non-negativity of the truncated
sum $\Hlocal(\sigma,\kappa)$ (Corollary~\ref{cor:Hlocal-pos});
critical divergence of order $(\log\kappa)^{2}$ on
$\sigma=\tfrac12$ (Lemma~\ref{lem:critical-div}); and
convergence to a finite limit for $\sigma>\tfrac12$
(Lemma~\ref{lem:subcrit-conv}).
These are unconditional mathematical statements.

This paper \emph{establishes numerically} two finite
computations from the stated parameters
(Observations~\ref{obs:cancellation}--\ref{obs:peak}):
the off-critical compensation between $\Hlocal$ and $\Hdual$
at $(\sigma,t,\kappa)=(0.3,0,1300)$, and the singular peak of
$\Hxi$ near $\gamma_1$ on the critical line.
Both numerical statements are finite computations from the
stated parameters and are labelled [numerical] in their
theorem headings.

This paper \emph{leaves open} the embedding of the singular
test function implicit in $\Hxi$ into the admissible
test-function class of the Weil functional
(Open Problem~\ref{op:weil-embedding}, \S\,\ref{sec:open}).
A positive answer would connect the curvature decomposition
to the Weil positivity framework; whether such a connection
extends to a full Weil positivity statement on the relevant
admissible class is a separate and deeper question.

% ============================================================
% §2 The Curvature Field
% ============================================================
\section{The Curvature Field}
\label{sec:field}

The Hadamard product representation of the entire function
$\xi(s)$, together with logarithmic differentiation, yields
\begin{equation}\label{eq:Hxi}
\Hxi(\sigma,t)
\;:=\;\frac{d^2}{d\sigma^2}\log\lvert\xi(\sigma+it)\rvert^2
\;=\;2\,\Re(\log\xi)''(s)
\;=\;-2\,\Re\sum_{\rho}\frac{1}{(s-\rho)^2},
\end{equation}
where the sum runs over the non-trivial zeros $\rho$ of
$\zeta(s)$ and $s=\sigma+it$.
The identity
$\partial_\sigma^2\log\lvert\xi(\sigma+it)\rvert^2
=2\,\Re[(\xi'/\xi)'(\sigma+it)]$
is explicit in
Gold\v{s}tejn--Grigutis~\cite{GoldstejnGrigutis24},
who study the sign of $\Re(\xi'/\xi)$ on horizontal lines;
our paper uses the same identity to decompose the
$\sigma$-Hessian into local and dual parts.
The quantity $\Hxi$ is the spectral (zero) side of an
explicit-formula identity at second-derivative level: formally
choosing a test function --- at the level of singular or
distributional test functions --- whose Mellin transform
produces the kernel $(s-\rho)^{-2}$ recovers $\Hxi$ as the
spectral output (rigorous treatment:
Open Problem~\ref{op:weil-embedding}).

The local (prime and archimedean) side is constructed
separately, in the adelic framework of Tate's thesis, as a
sum of per-place log-Mellin-norm contributions.
At the archimedean place,
\begin{equation}\label{eq:psi1-def}
\psi_1(\sigma)
\;:=\;\tfrac{1}{4}\,\frac{d^2}{d\sigma^2}
\log\lvert\Gamma(\sigma/2)\rvert^2.
\end{equation}
For real $\sigma$ this coincides with
$\tfrac{1}{8}\psi^{(1)}(\sigma/2)$, where $\psi^{(1)}$
denotes the trigamma function; one has
$\frac{d^2}{d\sigma^2}\log\lvert\Gamma(\sigma/2)\rvert^2
=\tfrac{1}{2}\psi^{(1)}(\sigma/2)$ by the chain rule, so
$\psi_1(\sigma)=\tfrac{1}{4}\cdot\tfrac{1}{2}
\psi^{(1)}(\sigma/2)=\tfrac{1}{8}\psi^{(1)}(\sigma/2)$.
At each finite place $p$, with $M_s$ the local Mellin
transform and $f_p$ the standard Tate test function,
\begin{equation}\label{eq:Vp-def}
V_p(\sigma)\;:=\;\frac{d^2}{d\sigma^2}
\log\lVert M_s f_p\rVert^{2}_{L^{2}(\mathbb{Q}_p^{\times})}.
\end{equation}
Explicitly, for the unramified local factor,
$\lVert M_s f_p\rVert^{2}=(1-p^{-1})(1-p^{-2\sigma})^{-1}$,
so that
\begin{equation}\label{eq:Vp}
V_p(\sigma)
\;=\;4(\log p)^{2}\,
\frac{p^{-2\sigma}}{(1-p^{-2\sigma})^{2}}.
\end{equation}
The truncated local curvature at cutoff $\kappa$ is
\begin{equation}\label{eq:Hlocal}
\Hlocal(\sigma,\kappa)
\;:=\;\psi_1(\sigma)\;+\;\sum_{p\le\kappa}V_p(\sigma),
\end{equation}
and the bookkeeping remainder is
\begin{equation}\label{eq:Hdual}
\Hdual(\sigma,t,\kappa)
\;:=\;\Hxi(\sigma,t)\;-\;\Hlocal(\sigma,\kappa).
\end{equation}
For $\sigma>\tfrac12$, $\Hlocal(\sigma,\kappa)$ converges as
$\kappa\to\infty$ (Lemma~\ref{lem:subcrit-conv}) and
$\Hdual(\sigma,t,\kappa)$ has a well-defined limit in
$\kappa$.
At $\sigma=\tfrac12$, $\Hlocal$ diverges
(Lemma~\ref{lem:critical-div}); the decomposition
$\Hxi=\Hlocal+\Hdual$ is then interpreted at every finite
$\kappa$, and $\Hdual$ is treated as a bookkeeping object
without spectral interpretation.

The decomposition $\Hxi=\Hlocal+\Hdual$ holds at every finite
cutoff $\kappa$ by the definition of $\Hdual$
in~\eqref{eq:Hdual}.
Read as a formal identity, it is the second-derivative
specialisation of the truncated explicit formula to the
singular test function whose Mellin side produces
$(s-\rho)^{-2}$; rigorous treatment of that test function is
the subject of Open Problem~\ref{op:weil-embedding}.
The identity holds by definition at every finite
cutoff~$\kappa$; the interpretation as a curvature-form
explicit-formula reorganisation is formal.

% ============================================================
% §3 Local Positivity
% ============================================================
\section{Local Positivity}
\label{sec:positivity}

\begin{observation}[Local positivity at finite places; proved]
\label{obs:Vp-pos}
For every prime $p$ and every $\sigma>0$,
\[
V_p(\sigma)\;\ge\;0.
\]
\end{observation}

\emph{Proof.}
From the explicit expression~\eqref{eq:Vp},
\[
V_p(\sigma)\;=\;4(\log p)^{2}\,
\frac{p^{-2\sigma}}{(1-p^{-2\sigma})^{2}}
\;=\;\Bigl(\frac{2(\log p)\,p^{-\sigma}}{1-p^{-2\sigma}}
\Bigr)^{2}\;\ge\;0.
\]
This is immediate: $V_p(\sigma)$ is a perfect square for every
$\sigma>0$, a special case of the classical log-convexity
of Dirichlet series with non-negative coefficients.
The content of the observation is the identification of
$V_p$ as the $\sigma$-Hessian of $\log\lVert M_s f_p\rVert^2$
in Tate's local framework~\cite{Tate1967}.
\hfill$\square$

\medskip
The archimedean contribution $\psi_1(\sigma)\ge 0$ is
classical, following from the log-convexity of the Gamma
function (Bohr--Mollerup theorem; cf.\ Artin~\cite{Artin}).
Combining the two place-wise positivities yields the
following.

\begin{corollary}[Truncated local curvature is non-negative;
                  proved]
\label{cor:Hlocal-pos}
For every $\sigma>0$ and every $\kappa\ge 2$,
\[
\Hlocal(\sigma,\kappa)\;\ge\;0.
\]
\end{corollary}

\medskip
Place-by-place positivity does not by itself approach the
Riemann Hypothesis (cf.\ \S\,\ref{sec:li-comparison}); its
role in this paper is to make the prime-side mass of $\Hxi$
explicit, so that the cancellation between $\Hlocal$ and the
bookkeeping remainder $\Hdual$ becomes visible at every finite
cutoff.
The numerical observations of \S\,\ref{sec:numerical}
quantify this cancellation off the critical line, where it is
near-total, and on the critical line, where $\Hxi$ exhibits
singular peaks near the non-trivial zeros.

% ============================================================
% §4 Divergence and Convergence
% ============================================================
\section{Divergence and Convergence}
\label{sec:divergence}

The behaviour of the truncated local curvature as
$\kappa\to\infty$ separates the half-plane into two regimes.
On the critical line, $\Hlocal$ diverges with order
$(\log\kappa)^{2}$; for every $\sigma>\tfrac12$, $\Hlocal$
tends to a finite limit.

\begin{lemma}[Critical divergence; proved]
\label{lem:critical-div}
As $\kappa\to\infty$,
\[
\Hlocal(\tfrac12,\kappa)
\;\asymp\;(\log\kappa)^{2}.
\]
\end{lemma}

\emph{Proof.}
From~\eqref{eq:Vp},
$V_p(\tfrac12)=4(\log p)^{2}p^{-1}(1-p^{-1})^{-2}$, so for
large $p$,
\[
V_p(\tfrac12)\;=\;4(\log p)^{2}\,p^{-1}\,(1+o(1)).
\]
Partial summation against the prime number theorem
$\pi(x)\sim x/\log x$ yields the standard Mertens-type
estimate
\[
\sum_{p\le\kappa}\frac{(\log p)^{2}}{p}
\;=\;\tfrac{1}{2}(\log\kappa)^{2}\;+\;O(\log\kappa)
\qquad(\kappa\to\infty)
\]
(Montgomery--Vaughan~\cite{MontgomeryVaughan}, \S\,2.2),
hence
$\sum_{p\le\kappa}V_p(\tfrac12)\asymp(\log\kappa)^{2}$.
The archimedean contribution $\psi_1(\tfrac12)$ is bounded,
so $\Hlocal(\tfrac12,\kappa)\asymp(\log\kappa)^{2}$.
Moreover,
\[
V_p(\tfrac12)
=4\frac{(\log p)^2}{p}+O\!\left(\frac{(\log p)^2}{p^2}\right),
\]
and the error term is summable over primes. Hence the same
Mertens-type estimate gives the sharper asymptotic
\[
\Hlocal(\tfrac12,\kappa)
=2(\log\kappa)^2+O(\log\kappa),
\]
with the bounded archimedean term absorbed into the error.
In particular, $\Hlocal(\tfrac12,\kappa)\asymp(\log\kappa)^2$.
\hfill$\square$

\begin{lemma}[Subcritical convergence; proved]
\label{lem:subcrit-conv}
For every $\sigma>\tfrac12$,
\[
\Hlocal(\sigma,\kappa)
\;\to\;\Hlocal(\sigma,\infty)\;<\;\infty
\qquad(\kappa\to\infty).
\]
\end{lemma}

\emph{Proof.}
From~\eqref{eq:Vp},
$V_p(\sigma)=O((\log p)^{2}p^{-2\sigma})$.
For $\sigma>\tfrac12$ one has $2\sigma>1$, hence the series
$\sum_{p}(\log p)^{2}p^{-2\sigma}$ converges absolutely, by
comparison with $\sum_p p^{-2\sigma+\delta}$ for any
$\delta\in(0,2\sigma-1)$.
The archimedean contribution $\psi_1(\sigma)$ is finite for
$\sigma>0$.
The result follows.
\hfill$\square$

\medskip
The critical line is therefore the threshold separating
divergence from convergence of the local curvature sum.
For completeness, the same comparison, again
by the prime number theorem, gives divergence for
$0<\sigma<\tfrac12$, since then
$\sum_p(\log p)^2p^{-2\sigma}$ diverges. Thus
$\sigma=\tfrac12$ is the boundary between convergence to the
right and divergence on and to the left of the critical line.

Matiyasevich, Saidak, and Zvengrowski~\cite{MSZ14}
identify $\sigma=\tfrac12$ as the conjectural boundary
of horizontal monotonicity of $\lvert\zeta(\sigma+it)\rvert$
via the first $\sigma$-derivative
(see also Sondow--Dumitrescu~\cite{SondowDumitrescu10}
for a related monotonicity reformulation of~RH).
Lemmas~\ref{lem:critical-div}
and~\ref{lem:subcrit-conv} establish the same critical line
as the divergence/convergence threshold of
$\Hlocal(\sigma,\kappa)$ via the second
$\sigma$-derivative; the two results are complementary.

% ============================================================
% §5 Numerical Illustration
% ============================================================
\section{Numerical Illustration}
\label{sec:numerical}

The two observations of this section quantify the
cancellation between $\Hlocal$ and $\Hdual$ at a
representative off-critical point and the singular peak of
$\Hxi$ near a non-trivial zero on the critical line.

\begin{observation}[Off-critical compensation; numerical]
\label{obs:cancellation}
At $(\sigma,t)=(0.3,0)$ with $\kappa=1300$,
\[
\Hlocal\;\approx\;915,
\qquad
\Hxi\;\approx\;0.09.
\]
The bookkeeping remainder $\Hdual$ therefore cancels $\Hlocal$
to better than $99.99\%$.
\end{observation}

\begin{observation}[Sampled values near zeros; numerical]
\label{obs:peak}
Take the rounded ordinates $t_1=14.13$, at distance
approximately $5\times 10^{-3}$ from the first non-trivial
zero $\gamma_1\approx 14.1347$, and $t_2=21.02$, at distance
approximately $2\times 10^{-3}$ from the second non-trivial
zero $\gamma_2\approx 21.0220$.
Direct evaluation of
$\Hxi(\sigma,t):=\partial_\sigma^{2}\log\lvert\xi(\sigma+it)\rvert^{2}$
at these ordinates gives
\[
\Hxi(0.3,t_1)\;\approx\;-49.78,
\qquad
\Hxi(\tfrac12,t_1)\;\approx\;+89\,577.92,
\]
and on the critical line near the second zero,
$\Hxi(\tfrac12,t_2)\approx +480\,754.91$.
\end{observation}

The spike at $\sigma=\tfrac12$ arises because the kernel
$2\Re(s-\rho)^{-2}$ in~\eqref{eq:Hxi} has a double pole at
$s=\rho$; near a zero this term dominates the spectral side,
producing the large values observed.
The larger value at $t_2$ in
Observation~\ref{obs:peak} reflects the smaller sampling
distance $\lvert t_2-\gamma_2\rvert\approx 2\times 10^{-3}$
to $\gamma_2$, compared with
$\lvert t_1-\gamma_1\rvert\approx 5\times 10^{-3}$ to
$\gamma_1$, together with the $\lvert s-\rho\rvert^{-2}$
singularity of the kernel.
The functional-equation symmetry
$\Hxi(\sigma,t)=\Hxi(1-\sigma,t)$, which follows from
$\xi(s)=\xi(1-s)$, holds to machine precision in all numerical
evaluations performed.

Together, Observations~\ref{obs:cancellation}
and~\ref{obs:peak} quantify the cancellation between the
place-by-place non-negative $\Hlocal$
(Observation~\ref{obs:Vp-pos}, Corollary~\ref{cor:Hlocal-pos})
and the bookkeeping remainder $\Hdual$: numerically, the
cancellation reaches better than $99.99\%$ off the critical
line, while on the critical line $\Hxi$ takes large positive
values when $(\sigma,t)$ is sampled close to a non-trivial
zero, consistent with the singular kernel $(s-\rho)^{-2}$.
The relationship is observational at finite cutoff and is
not asserted as a theorem.

% ============================================================
% §6 Rankin-Selberg Analogue
% ============================================================
\section{Rankin--Selberg Analogue}
\label{sec:rs}

The divergence at $\sigma=\tfrac12$ established in
Lemma~\ref{lem:critical-div} is not specific to $\zeta(s)$.
This section is informational.
Under the standard Rankin--Selberg prime-number/Tauberian
framework for a cuspidal automorphic representation~$\pi$,
and after removing ramified primes and separating prime-power
terms, one expects the weighted prime sum
\[
\sum_{p\le\kappa}\lvert\lambda_{\pi}(p)\rvert^{2}\,
\frac{(\log p)^{2}}{p}
\]
to have order $(\log\kappa)^{2}$, with a positive leading
constant depending on the normalisation of the
Rankin--Selberg residue
(cf.\ Kowalski--Michel~\cite{KowalskiMichel} for related
prime sums in the Rankin--Selberg context).
This is not used in any proof of the present paper.

The Rankin--Selberg $L$-function extension of Li's
criterion has been studied by
Smajlovi\'c~\cite{Smajlovic10}
and Od\v{z}ak--Smajlovi\'c~\cite{OdzakSmajlovic10}.
The analogy in the present paper is formal and
motivational; no Li-criterion proof is claimed here.

% ============================================================
% §7 Connection to Li Coefficients
% ============================================================
\section{Connection to Li Coefficients}
\label{sec:li-comparison}

The Li coefficients $\lambda_n$, defined via
\[
g(z)\;:=\;\log\xi\!\left(\frac{1}{1-z}\right),
\qquad
g'(z)\;=\;\sum_{n\ge 1}\lambda_n\,z^{n-1},
\]
satisfy Li's criterion~\cite{Li}:
\[
\mathrm{RH}\;\Longleftrightarrow\;
\lambda_n\ge 0\quad\text{for all }n\ge 1.
\]
Differentiating once more,
\begin{equation}\label{eq:gpp}
g''(z)\;=\;(\log\xi)''\!\left(\frac{1}{1-z}\right)(1-z)^{-4}
\;+\;2(\log\xi)'\!\left(\frac{1}{1-z}\right)(1-z)^{-3},
\end{equation}
so that the term $(\log\xi)''(s)$ appearing in~\eqref{eq:gpp}
is exactly the core of the curvature field
$\Hxi(s)=2\Re(\log\xi)''(s)$.
The present curvature quantities are therefore structurally
close to Li-type positivity criteria: both are built from
logarithmic derivatives of $\xi$.
The difference is that $\Hxi$ is a continuous pointwise
second-derivative field, whereas the Li coefficients package
the same logarithmic data into a discrete sequence after the
M\"obius reparametrisation around $s=1$.

\begin{remark}
The local positivity $V_p(\sigma)\ge 0$ at each finite place
established in Observation~\ref{obs:Vp-pos} does \emph{not} imply
$\lambda_n\ge 0$.
Li positivity is a global statement mixing all places, the
functional equation, the complete zero geometry, and a
specific test-function structure on the disc parametrisation.
Place-by-place positivity is a compatible local input but not
sufficient for Li positivity, and the gap between the two is
of Riemann-Hypothesis strength.
\end{remark}

% ============================================================
% §8 Open Problems
% ============================================================
\section{Open Problems}
\label{sec:open}

The central open problem arising from the curvature
decomposition of this paper is the following.

\begin{openproblem}[Embedding into the Weil framework]\label{op:weil-embedding}
Does there exist a global packaging --- a summation
procedure, test-function class, or regularisation --- that
transfers the local positivity
\[
V_p(\sigma)\ge 0\;(\forall\,p\,\text{prime}),
\qquad
\psi_1(\sigma)\ge 0,
\]
together with the critical divergence
$\Hlocal(\tfrac12,\kappa)\to\infty$, into a global
Li- or Weil-type positivity statement?

More precisely: can the singular test function implicit in
$\Hxi$ --- whose Mellin side produces the kernel
$(s-\rho)^{-2}$ --- be embedded in, or approximated by, the
class $\Sad$ of admissible test functions for which Weil
positivity $W(f*\check f)\ge 0$ is equivalent to the Riemann
Hypothesis?
\end{openproblem}

We note that Weil positivity is highly sensitive to the
choice of test-function class and regularisation.
Bombieri~\cite{Bombieri2000} studies Weil's quadratic
functional on support-restricted $L^2$-spaces and its
finite-dimensional truncations, while
Lagarias~\cite{Lagarias2007} formulates the Li/Weil
framework on test spaces with explicit analyticity and
decay restrictions. These results show that one cannot
pass from a formal explicit-formula kernel to a positivity
statement without controlling the admissible class and
the limiting process. The kernel implicit in the present
paper, whose Mellin side produces $(s-\rho)^{-2}$, is more
singular than the first-order distributions appearing in
the classical explicit formula. The problem of embedding
or approximating it inside an admissible Weil framework
is therefore a genuine regularisation problem, not a
formal consequence of local positivity.

A positive answer to Open Problem~\ref{op:weil-embedding}
would connect the curvature decomposition to the Weil
positivity framework
(Weil~\cite{Weil1952}; Lagarias~\cite{Lagarias};
Connes--Consani~\cite{ConnesConsani}).
By Weil's criterion
(Weil~\cite{Weil1952}; Bombieri--Lagarias~\cite{BombieriLagarias}),
positivity of the Weil quadratic functional on the full relevant
admissible class is equivalent to the Riemann Hypothesis.
Positivity for the particular family arising from the present
construction would therefore be only a family-level positivity
statement, not by itself an RH criterion.
The present Open Problem is therefore the prior step of
fitting $\Hlocal(\sigma,\kappa)$ into the admissible
framework, and is not the positivity step itself; the latter
remains a separate and deeper open question, and is not
addressed here.

% ============================================================
% §9 Non-Circularity
% ============================================================
\section*{Non-Circularity}
\label{sec:nocirc}

We document explicitly that the results of this paper do not
rely on the objects they are directed toward.

\textbf{No RH as input.}
The Riemann Hypothesis is not assumed at any point.
The non-trivial zeros $\rho$ enter the curvature field $\Hxi$
through the Hadamard product representation of $\xi$, which
is unconditional; no property of the location of the zeros is
used in the formal core
\S\S\,\ref{sec:field}--\ref{sec:divergence}.

\textbf{No GUE, Montgomery, or Hilbert--P\'olya hypothesis.}
No correlation statistic of the zeros and no postulated
spectral realisation of the zeros is assumed or used.
The terminology of ``curvature'' refers solely to the
second-derivative origin of $V_p(\sigma)$ in
$\log\lvert\xi(\sigma+it)\rvert^{2}$ and is not used as a
geometric or operator-theoretic input.

\textbf{No Li positivity as input.}
The Li criterion of \S\,\ref{sec:li-comparison} is cited only
as a comparison object; positivity of the coefficients
$\lambda_n$ is not assumed and not used in any proof of this
paper.
The Remark in \S\,\ref{sec:li-comparison} explicitly states
that local positivity does not imply Li positivity.

\textbf{No Weil positivity as input.}
The Weil functional $W(f*\check f)$ and the admissible class
$\Sad$ enter only in the formulation of
Open Problem~\ref{op:weil-embedding}.
No positivity statement on $W$ is assumed in any of
Observation~\ref{obs:Vp-pos}, Lemmas~\ref{lem:critical-div},
\ref{lem:subcrit-conv}, or in
Corollary~\ref{cor:Hlocal-pos}.

\textbf{Inputs from classical analytic number theory.}
The formal core of the paper uses only standard inputs:
the Hadamard product for $\xi$, log-convexity of the Gamma
function (Bohr--Mollerup), the prime number theorem
$\pi(x)\sim x/\log x$, and elementary partial summation.
The Rankin--Selberg discussion in \S\,\ref{sec:rs} relies
additionally on the simple pole of
$L(s,\pi\times\widetilde{\pi})$ at $s=1$ and a Tauberian
argument; this section is informational and not used in the
formal core.

% ============================================================
% Acknowledgments
% ============================================================
\section*{Acknowledgments}

The author thanks Jean-Fran\c{c}ois Burnol for confirming
that the prime-cutoff curvature formulation does not appear
in the literature known to him, and for suggesting connections
to the noncommutative-geometry programme of Connes and to the
Langlands framework.
The author also thanks the readers who provided feedback on
earlier versions of this note.

% ============================================================
% Call for Feedback
% ============================================================
\section*{Call for Feedback}

This paper is part of an ongoing open research initiative
toward the Riemann Hypothesis.
All papers in the series are freely available on Zenodo under
CC~BY~4.0, and the accompanying verification code is hosted
on GitHub.

{\sloppy
Zenodo (\url{https://zenodo.org/search?q=Ulrich+Tehrani})
and GitHub
(\url{https://github.com/utehrani/analysislab-nt}).
\par}

Topics of particular interest for diagnostic feedback on this
note include:
(i) whether the singular kernel $(s-\rho)^{-2}$ underlying
$\Hxi$ admits a natural regularisation or limit construction
in $\Sad$, in the sense of
Open Problem~\ref{op:weil-embedding}; and
(ii) whether the Rankin--Selberg analogue
(\S\,\ref{sec:rs}) admits a refined asymptotic with explicit
leading constant under standard conventions.

% ============================================================
% References
% ============================================================
\newpage
\begin{thebibliography}{99}

\bibitem{Connes}
A.~Connes,
``Trace formula in noncommutative geometry and the zeros of
the Riemann zeta function'',
Sel.\ Math.\ (N.S.) \textbf{5}(1), 29--106 (1999).

\bibitem{ConnesConsani}
A.~Connes and C.~Consani,
``Weil positivity and trace formula, the archimedean place'',
Sel.\ Math.\ (N.S.) \textbf{27}, 77 (2021).

\bibitem{Weil1952}
A.~Weil,
``Sur les `formules explicites' de la th\'eorie des nombres
premiers'',
Comm.\ S\'em.\ Math.\ Univ.\ Lund [Medd.\ Lunds Univ.\ Mat.\ Sem.],
Tome Suppl\'ementaire (1952), 252--265.

\bibitem{Lagarias}
J.~C.~Lagarias,
``On a positivity property of the Riemann $\xi$-function'',
Acta Arithm.\ \textbf{89}(3), 217--234 (1999).

\bibitem{BombieriLagarias}
E.~Bombieri and J.~C.~Lagarias,
``Complements to Li's criterion for the Riemann hypothesis'',
J. Number Theory \textbf{77}(2), 274--287 (1999).

\bibitem{Li}
X.-J.~Li,
``The positivity of a sequence of numbers and the Riemann
hypothesis'',
J. Number Theory \textbf{65}, 325--333 (1997).

\bibitem{Burnol}
J.-F.~Burnol,
``Sur les formules explicites I: analyse invariante'',
C.\ R.\ Acad.\ Sci.\ Paris S\'er.\ I Math.\ \textbf{331}(6),
423--428 (2000).

\bibitem{KowalskiMichel}
E.~Kowalski and P.~Michel,
``Zeros of families of automorphic $L$-functions close to 1'',
Pacific J.\ Math. \textbf{207}(2), 411--431 (2002).

\bibitem{Artin}
E.~Artin,
\emph{The Gamma Function},
Holt, Rinehart and Winston, New York, 1964
(translation by M.~Butler of \emph{Einf\"uhrung in die Theorie
der Gammafunktion}, Teubner, Leipzig, 1931).

\bibitem{MontgomeryVaughan}
H.~L.~Montgomery and R.~C.~Vaughan,
\emph{Multiplicative Number Theory I: Classical Theory},
Cambridge Studies in Advanced Mathematics \textbf{97},
Cambridge University Press, 2006.

\bibitem{Edwards}
H.~M.~Edwards,
\emph{Riemann's Zeta Function},
Academic Press, New York, 1974.

\bibitem{BerryKeating1999}
M.~V.~Berry and J.~P.~Keating,
``$H=xp$ and the Riemann zeros'',
in \emph{Supersymmetry and Trace Formulae},
eds.\ I.~V.~Lerner et al.,
Kluwer/Plenum, New York, 1999, pp.\ 355--368.

\bibitem{deBranges}
L.~de~Branges,
\emph{Hilbert Spaces of Entire Functions},
Prentice-Hall, Englewood Cliffs, NJ, 1968.

\bibitem{deBranges1986}
L.~de~Branges,
``The Riemann Hypothesis for Hilbert Spaces of Entire
Functions'',
\emph{Bull.\ Amer.\ Math.\ Soc.\ (N.S.)} \textbf{15}(1),
1--17 (1986).

\bibitem{Tate1967}
J.~Tate,
``Fourier Analysis in Number Fields and Hecke's
Zeta-Functions'',
in \emph{Algebraic Number Theory},
eds.\ J.~W.~S.~Cassels and A.~Fr\"ohlich,
Academic Press, London, 1967, pp.\ 305--347
(PhD thesis, Princeton University, 1950).

\bibitem{GoldstejnGrigutis24}
E.~Gold\v{s}tejn and A.~Grigutis,
``On a positivity property of the real part of the
logarithmic derivative of the Riemann $\xi$-function,''
\emph{J.\ Math.\ Inequal.} \textbf{18}(3), 829--845 (2024).
DOI: 10.7153/jmi-2024-18-45.

\bibitem{Stopple16}
J.~Stopple,
``Notes on $(\log\zeta(s))''$,''
\emph{Rocky Mountain J.\ Math.} \textbf{46}(5) (2016),
1701--1715; arXiv:1311.5465.

\bibitem{Bombieri2000}
E.~Bombieri,
``Remarks on Weil's quadratic functional in the theory
of prime numbers, I,''
\emph{Rend.\ Mat.\ Acc.\ Lincei}, Ser.~IX,
\textbf{11} (2000), 183--233.

\bibitem{Lagarias2007}
J.~C.~Lagarias,
``Li coefficients for automorphic $L$-functions,''
\emph{Ann.\ Inst.\ Fourier} \textbf{57} (2007),
1689--1740.

\bibitem{MSZ14}
Yu.~Matiyasevich, F.~Saidak, and P.~Zvengrowski,
``Horizontal monotonicity of the modulus of the
Riemann zeta function, $L$-functions, and related
functions,''
\emph{Acta Arith.} \textbf{166}(2) (2014), 189--200;
arXiv:1205.2773.

\bibitem{Haran90}
S.~Haran,
``Riesz potentials and explicit sums in arithmetic,''
\emph{Inventiones Math.} \textbf{101} (1990), 697--703.

\bibitem{Smajlovic10}
L.~Smajlovi\'c,
``On Li's criterion for the Riemann hypothesis for
the Selberg class,''
\emph{J.\ Number Theory} \textbf{130}(4) (2010),
828--851.

\bibitem{OdzakSmajlovic10}
A.~Od\v{z}ak and L.~Smajlovi\'c,
``On Li's coefficients for the Rankin--Selberg
$L$-functions,''
\emph{Ramanujan J.} \textbf{21}(3) (2010), 303--334.

\bibitem{SourmelidisSteuding24}
A.~Sourmelidis and J.~Steuding,
``Spirals of Riemann's zeta-function --- curvature,
denseness and universality,''
\emph{Math.\ Proc.\ Cambridge Philos.\ Soc.}
\textbf{176}(2) (2024), 325--338; arXiv:2306.00460.

\bibitem{FedosovaRowlettZhang20}
K.~Fedosova, J.~Rowlett, and G.~Zhang,
``Second variation of Selberg zeta functions and
curvature asymptotics,''
\emph{Ann.\ Global Anal.\ Geom.} \textbf{57} (2020),
23--60; arXiv:1709.03841.

\bibitem{SondowDumitrescu10}
J.~Sondow and C.~Dumitrescu,
``A monotonicity property of Riemann's xi function
and a reformulation of the Riemann hypothesis,''
\emph{Period.\ Math.\ Hungar.} \textbf{60}(1) (2010),
37--40; arXiv:1005.1104.

\end{thebibliography}

% ============================================================
% Footer -- normative format
% ============================================================
\enlargethispage{3\baselineskip}
\vspace*{\fill}
\begin{minipage}{\linewidth}
\rule{\linewidth}{0.4pt}\smallskip\\
\small\emph{Paper~1 v23 \textperiodcentered\ May~2026
\textperiodcentered\ Pauca sed matura}
\end{minipage}

\end{document}
